\documentclass{jsarticle}
\usepackage{amsmath}
\usepackage{amssymb}
\begin{document}
\title{Hilbert manifold and Seedist of Categories theorem \\
Surface of deprivation and One projective \\
surface released into All of dimension.}
\author{Masaaki Yamaguchi, Michihiro Yamaguchi, Tokiko Yamaguchi \\
(ichuga my is own selves \\
from out of skills and potential power included in.) \\
(ichuga my is uo roof teams, \\
and this skills create potential aquire of knowleges \\
from this safty powers.) \\
(Tokiko means to arrow of target ones and forever loves with family. \\
Toki call them to messure time to \\ 
relativty theory of special and general.) \\
Sachie Yamaguchi, Mutsuyo Wada,Nobuko Yamaguchi \\
(Yo is world and no relate of tsumu of Yo, \\
and this security of world \\
call her to prevent \\
world to safe them to other killers. \\
Moreover, this skills call them to own skill.) \\
(ichuga my achi as is meant with see of \\
next days of incident and event moreover, \\
see her to know the future events.)
(oku bon is my brother. books no relate \\
this to call him to name.
After all, He promised me to being \\
no related with Homo into other man.) \\
(and tokyo construct with "No" of being named \\
from build of hill and city trast.) \\
And add to this report with \\
this precious person call me to birth of my son \\
in present of acasic record.}
\date{}
\maketitle

山口雅旭は、山口家のみやびのあさひであり、夜と昼である。\\
山口道弘は、山口家の弘法大師さまの推算式であり、行いである。 \\
山口トキ子は、山口家の相対性理論の場の時間であり、酉のトキであり、弓道である。 \\
山口幸恵は、山口家の海の幸であり、幸せの恵みである。 \\
山口順子は、山口家の伝令の神の子であり、ルールを守る人である。 \\
和田睦世は、和田家の技の世を守る人であり、神様の子である。 \\
アカシックレコードの子は、情報空間のビデオを見ている子であり、 \\
その場に滞在している子である。 \\
ユウくんは、優しさの奉仕、優秀な頭脳、友だち、夕暮れ、ゆうパックである。\\
モネは、クロード・モネのカミユさんの日本画から由来している名前である。 \\
綺麗さが際立っていて、私が、扇の絵に好いて、モネに命名した名前である。\\
全員、侍から命名されている。
山口知世子おばさんは、山口家の大善の予知力を司る役割を受け持っている。 \\
山口香織さんは、山口家の瞑想のための線香の場の雰囲気を醸し出す、 \\
役割を担っている大善である。それで、大それている名前らしい。 \\
そのために、つかみ所のない人になっているらしい。 \\
瞑想法で読まれないとは、すごすぎる。 \\
山口雅旭は、既に述べた通りの、山口彩さんが、私の家系に嫁ぎに来ると、 \\
私の朝になってくれる賁臨であるらしい。
Bubdaの瞑想法が、それらしい。
浄土真宗の親鸞和尚と達磨大師、曹洞宗が、戒律を緩めて、如伴を受け持ってもいいと、\\
助言してくれていて、そのことを、示唆した男性が、\\
五木寛之氏を、私は、好いているが、ダメと言われていたが、 \\
嶋本淑子さんだけは、許すと言ってくれた、\\
恥ずかしいのか、年賀状を棄却された人に、 \\
蔵王病院であった。
仏壇に親鸞和尚を飾っているらしく、仏壇は、\\
仏の棺であるが、ブツを断するらしい。
時と場合のことを言っている。
握手を交わすとは、読んでいる、読んでいない、仏教の
断りを交わすかどうかであり、 \\
奥が深い哲学である。\\
浄土真宗は、念仏を唱えて、自身が空になり、仏と一体化することであり、\\
悪人正機とは、以前の苦悪を、念仏を唱えることによって、自身が空になり、\\
瞑想法の修行により、周りに読まれる訓練によって、していた苦悪を浄化する、\\
そのための念仏であり、如伴とは違う意味らしく、\\
真言宗の霊山寺と同じ役割をされる念仏であるので、\\
意見が合わずに、握手を交わさなかったらしい。\\
山口裕子さんは、融通性が効く優しさの奉仕であり、懐の深い哲学である。\\
山口勝人さんは、人の上に立つ人である。\\
山口雅代さんは、雅な題目を掲げる人である。礼節と手本である。\\
塩飽弘さんは、潮汐力を軽減する人であり、弘法大師の現身である。\\
塩飽由美子さんは、潮汐力を軽減する人であり、弓道の元である。\\
苫米地英人博士は、占いを維持しているアメリカの地の人であり、\\
ダンディーであり、ラジャループと紳士を併せ持った人である。
石川博先生は、不動の水星であり、博士の本元である。\\
河相和昭先生は、水星と木星の相を併せ持った人であり、日本の昭和の初めの \\
初心を持った人である。 \\
脇本和博先生は、月の力の結晶のアカシックレコードの一部と、\\
日本の博士の本元である。\\
河相葉子先生は、水星と木星の相を合わせていて、植物と多様な成分を分解される \\
先生であり、心理士としての、分析力を発揮される人である。 \\
馬屋原健先生は、魔法の薬と、その対象の健康との戦いの人という立場の \\
人である。 \\
馬屋原壮六先生は、魔法の薬と、優しさと厳しさの荘厳の人である。 \\
奥田伴枝薬剤師さんは、大善を守り、魔法の亥年としてと、枝葉末節 \\
としての、植物の分岐からくる、生命の神秘の人である。 \\
平田渉美さんは、平常心の状態での綺麗さで歩く人である。\\
藤田靖子さんは、ドルフィンクラブの日本の元である。\\
上さんは、私のど真ん中のタイプである。 \\
杜裕子さんは、寺院の鳥居であり、優しさの奉仕である。 \\
和田絵理香さんは、所属している役の技に長けている人であり、\\
速読の能力で、物事の断りを瞑想法の場によって、\\
読む能力に長けている人である。\\
河原真澄さんは、京都の鴨川沿いでの、奉仕の賁臨である人である。\\
要するに、水星の魔法である。\\
津村雅美さんは、漢方薬の賁臨である。食事の元でもある。食神の人でもある。\\
斎藤一人先生は、抜刀齊の山口一の人であり、千人斬りの人と、\\
火星の酉の人である。\\
石岡芳隆先生は、不動の剣道の、仏教でのお香の場を作り出す人である。\\
栗田昌祐先生は、酉の木星を守る人であり、朝と昼の優しさと真逆の夜、\\
全ての時間の役割を担っている優しさである。\\
川島隆太教授は、しだれ桜の水星であり、心太の元である。\\
山口信先生は、全干支の出入り口の役割を、周囲の人が信徳として信じられる \\
役割を発揮する人である。\\
立花隆先生は、木星の立花である陽の健康を表すひとである。陰と陽の陰陽学 \\
である。\\
安岡正篤先生は、安心していられる健康体であり、篤姫の陰の立場の人である。\\
細木数子先生は、繊細なこころの木星であり、数学の女王である。\\
細木かおり先生は、仏教の人である。\\
棗田和美薬剤師さんは、柑橘の薬であり、まぼらほの魔法の人であり、日本の女性である。
嶋本淑子さんだけは、紳士淑女であるのを、関係を持つと、アカシックレコードの一部を \\
もらえた人である。\\
渡辺由美さんは、水星の推算式であり、弓道の優雅さである。\\
諸藤真理さんは、ど真ん中の日本であり、世の中の道理の真理である。\\
中谷さんは、繊細なこころで、物事を中てる人であり、溝口さんは、水星の努力家であり、\\
前原一之さんは、魔法の薬の水星であり、推算式で一筋である人、\\
徳山さんは、陰徳と陽徳の水星であり、徳永さんは、徳の永遠の使命であり、\\
松浦まどかさんは、魔法の水星と木星の相である裏と表の力の結晶のアカシックレコードであり、\\
松枝さんは、魔法の水星と木星の能力の源である、枝葉末節の能力であり、\\
水口さんは、水星の能力の投射と関与の能力であり、\\
田中さんは、ルールを守っていることによる、物事を中てる人であり、\\
中田由香里さんは、物事を中てるために、ルールを守り、弓道の優雅さで、御先祖様を大切にする\\
人であり、松本さんは、アカシックレコードの一部である魔法の木星であり、\\
掛ヶ谷さんは、優しさと計算力を授かる人であり、時繁剛さんは、時間の魅力を自身の強さで\\
味わう人であり、長田さんは、中てる殺菌の人であり、ルールを守るための使命であり、\\
若山さんは、水星の初代運の人であり、\\
名前の通りには、人生は過ごせないとも、多種多様な人生がある。\\
立春は春に至り、春分の日は、春のど真ん中であり、仕切りの力があり、\\
立夏は、夏が経ち、夏至は、夏に至り、立秋は、秋が経ち、秋分は、木星の相であり、\\
秋のど真ん中であり、立冬は、冬が立ち、冬至は、冬のど真ん中であり、大晦日は、\\
大善の最後の仏の優しさで、仕切りになり、元旦は、1年の始まりの日であり、\\
八十八夜は、仏教の如伴を許される日であり、医学と易学、暦は、世の中の道理を\\
知るための学問になっている。 \\
生理周期は、この八十八夜の3ヶ月での生理と月経のためのリンゴ酸回路とクエン酸サイクル\\
と同じ原理で、エストロゲンに精子を付着させるタイミングのための準備期間であり、\\
3週間間隔らしく、体外受精がヒントになっている。\\
食べ物の消化が、どの成分がどの臓器で消化されて、吸収されるかが、\\
このエストロゲンに精子を付着させるタイミングのヒントになっている。\\
私の精子は、塩素の匂いがして、水の中に塩素が入っている匂いになっている。\\
相当、私の子の推算式が期待大である。\\
大坂くんは、土星と水星の大善であり、広島大学医学部の看護師さんの、\\
間所さんは、もんどころの関所であり、関税自主権としての政府の関所であり、\\
アカシックレコードの門としての、関所である。\\
未光ひとみさんは、未だ未にあらずであり、希有でのつかみ所の無い人であり、目である。\\
横田則夫先生は、横断歩道を守り、水星の陽の人である。真逆は、オリオン後に恋するである。\\
後藤謙一朗先生は、日本の後の代の守りである。語ののれんの侍である。\\
辻先生は、武士である。\\
梶山先生は、土星と水星の尾である。\\
萬谷智之先生は、速読法の瞑想法であり、智性の易である。\\
中村先生は、木星の寸により、物事を中てる先生である。\\
渡辺千種先生は、水星と木星の相で、多種多様な能力を備えた人である。\\
アカシックレコードの一部である魔法を扱う人である。\\
福馬先生は、幸福を約束する魔法の午の人である。\\
山口香織先生が、鮎沢香織先生になっていて、鮎沢香織先生は、イタリアの鮎料理であり、水星の占い\\
であり、お香の瞑想法の場を作り出す人である。\\
鮎沢吉孝は、水星の占いの土星を切り、そして、それを手にする吉凶の人である。\\
鮎沢康江お姉さんは、水星の占いであり、川端康成氏の名前と、金星の陰の速読の人である。\\
鮎沢良彦は、水星の占いであり、武士である。\\
鮎沢知代は、水星の占いであり、全てを知る人である。\\
美代子おばあちゃんは、美しさを代々受け持つ人であり、大善の黄泉さまである。\\

食べ物で、キャベツは、kyabetu$\to$tue ab i ki$\to$mars day sunday ni uete hiru moku ni toru\\
$\to$日曜の昼に植えて、3月から4月にかけて、8月の木曜に収穫する。\\
レタスは、retasu$\to$u no amasa no hi ni ueru torunoha 3kagetugo\\
$\to$3月の3週間の火曜に植えて、6月の3週間後の木曜に収穫する。\\
スイカは、suika$\to$kayou no i ni su de uete 4gatu ni ueru 3,4kagetugo ni shuukaku suru.\\
緑黄色野菜は、3月から4月にかけて植えて、3，4ヶ月後に収穫する。\\
果物は、4月に植えて、7，8月に収穫する。\\
イチゴは、1月に植えて、5月に収穫する。\\
メロンは、スイカと同じく、4月に植えて、meron$\to$ ataru en$\to$9gatu ni shuukaku suru.\\
リンゴは、ringo$\to$ 3gatu no sora ni uete naru rin ga nakunarunoha 9gatu ni shuukaku suru.\\
in no tabemono ha you no kankyou de dekite, you no tabemonoha in no kankyou de dekiru.\\
onmyousoutairon de aru.\\
meron ha atuitokoro de dekite ringo ha samui tokoro de dekiru. suika ha in de you no 
tokorode dekiru. retasu ha in de youno tokorode dekiru. kyabetu mo in de youno tokorode dekiru.
korerayori dono hougaku de tabemono ga dekiruka wo onmyousouutaironn ha nobeteiru.
果物は、火曜に植えて、金曜に収穫する。
野菜は、火曜に植えて、木曜に収穫する。
お芋さんは、火曜に植えて、金曜に収穫する。お芋さんは、4月に植えて、11月に収穫する。
紅天使さんが、お芋さんの代表である。
 
\section{Hilbert manifold in Mebius space \\
        this element of Zeta function on integrate of fields}

 Hilbert manifold equals with Von Neumann manifold, and this fields is concluded with Glassman manifold,
the manifold is dualty of twister into created of surface built.
These fields is used for relativity theorem and quantum physics.
$$ ||ds^2|| = \lim_{x \to \infty} [\delta(x) \int \int \int \pi \left( \sum_{k=0}^{\infty}
		{{}^n\sqrt{p},x \over n} \right)
^{1 \over 2}d \tau]^{\mu\nu} $$
  This norm is component of fields in Hilbert manifold of space theorem rebuilt with
Mebius space in Gamma Function on Beta function fill into power of rout fundamental
group.
And this result with $AdS_5$ space time in Quantum caos of Minkowsky of manifold in
abel manifold.
Gravity of metric in non-commutative equation is Global differential equation
conbuilt with Kaluza-Klein space. Therefore this mechanism is $T^{\mu\nu}$ tensor
is equal with $R^{\mu\nu}$ tensor. And This moreover inspect with laplace operator
in stimulate with sign of differential operator.
Minus of zone in Add position of manifold is Volume of laplace equation rebuilt with
Gamma function equal with summative of manifold in Global differential equation,
this result with setminus of zone of add summative of manifold, and construct with
locality theorem straight with fundamental group in world line of surface, this
power is boson and fermison of cone in hyper function.

$$ V(\tau) = [f(x),g(x)] \times [f^{-1}(x), h(x)] $$
$$ \Gamma(p,q) = \int e^{-x} x^{1 - t}dx $$
$$ = \beta(p,q) $$
$$ = \pi (f(\chi,x),x) $$
$$ ||ds^2|| = \mathcal{O}(x)[(f(x) \circ g(x))^{\mu\nu}]dx^{\mu}dx^{\nu} $$
$$ = \lim_{x \to \infty} \sum_{k=0}^{\infty}a_k f^k $$
$$ G^{\mu\nu} = {\partial \over \partial f}\int [f(x)^{\mu\nu} \circ G(x)^{\mu\nu}dx^{\mu}dx^{\nu}]^{\mu\nu}dm $$
$$ = g_{\mu\nu}(x)dx^{\mu}dx^{\nu} - f(x)^{\mu\nu}dx^{\mu}dx^{\nu} $$
$$ [i \pi(\chi,x),f(x)] = i \pi f(x) - f(x) \pi(\chi,x) $$
$$ T^{\mu\nu} = (\lim_{x \to \infty} \sum_{k=0}^{\infty}\int \int [V(\tau) \circ S^{\mu\nu}(\chi,x)]dm)
^{\mu\nu}dx^{\mu}dx^{\nu} $$
$$ G^{\mu\nu} = R^{\mu\nu}T^{\mu\nu} $$
$$ \begin{vmatrix} D^m & dx \\
dx & \sigma^m \end{vmatrix} \begin{vmatrix} \cos \theta & - \sin \theta \\
\sin \theta & \cos \theta \end{vmatrix} \begin{vmatrix} x \\
y \end{vmatrix} = \begin{pmatrix} 1 & 0 \\ 
0 & - 1 \end{pmatrix}^{1 \over 2} $$
$$ \sigma^m \begin{bmatrix} \delta(x) & - 1 \\
1 & \epsilon(x) \end{bmatrix}^{1 \over 2} = \begin{pmatrix} i & 0 \\
0 & - i \end{pmatrix} $$
$$ V(M) = {\partial \over \partial f} ({}^N \int [f \smallsetminus M]^{\oplus N})^{\mu\nu} dx^{\mu}dx^{\nu} $$
$$ V(M) = \pi(2 \int \sin^2 dx)\oplus {d \over df}F^M dx_m $$
$$ \lim_{x \to \infty}\sum_{k=0}^{\infty} a_k f^k = \int (F(V)dx_m)^{\mu\nu}dx^{\mu}dx^{\nu} $$
$$ \bigoplus_{k=0}^{\infty}[ f \smallsetminus g] = \vee(M \wedge N) $$
$$ \pi_1 (M) = e^{-f 2 \int \sin^2 x dm} + O(N^{-1}) $$
$$ = [i \pi(\chi,x),f(x)] $$
$$ M \circ f(x) = e^{-f \int \sin x \cos x dx_m} + \log (O(N^{-1})) $$

 Non-Symmetry space time.
 $$ {d \over dL}V(\tau) = {d \over df}\int \int_M {1 \over (x \log x)^2}dx_m + {d \over df}\int \int_M
 {1 \over (y \log y)^{1 \over 2}}dy_m $$
 $$ \epsilon S(\nu) = \square_v \cdot {\partial \over \partial \chi}({}^5 \sqrt{\wedge g^2})d\chi $$
 Differential Volume in $AdS_5$ graviton of fundamental rout of group.
 $$ \wedge (F^m_t)^{''} = {1 \over 12}g_{ij}^2$$
 Quarks of other dimension.
 $$ \pi(V_{\tau}) = {e^{-(\sqrt{\pi \over 16}\log x)}}^{\delta} \times {1 \over (x \log x)} $$
 Universe of rout, Volume in expanding space time.
 $$ {d \over dt}(g_{ij})^2 = {1 \over 24}(F^m_t)^2 $$
 $$ m^2 = 2 \pi T \left({26 - D_n \over 24}\right) $$
 This quarks of mass in relativity theorem, and fourth of universe in three manifold of one dimension surface,
 and also this integrate of dimension in conbuilt of quarks.
 $$ g_{ij} \wedge \pi(\nu_{\tau}) = e^{-2 \pi T |\psi|}[\eta_{\mu\nu} + \bar{h}_{\mu\nu}(x)]dx^{\mu}dx^{\nu} 
 + T^2 d \psi^2 $$
 Out of rout in $AdS_5$ space time.
 $$ ||ds^2|| = g_{ij}\wedge \pi(\nu_v) $$
 $AdS_5$ norm is fourth of universe of power in three manifold out of rout.


\section{Entropy}
   Already discover with this entropy of manifold is blackhole of energy, these entropy is mass of energy has
norm parameter. And string theorem $ T_{\mu\nu}$  belong to $ \rho $ energy, this equation is same Kaluze-Klein theorem.


 $$ \nabla \phi^2 = 8 \pi G({p \over c^3} + {V \over S})$$

 $$  y = x $$
 $$  y = (\nabla \phi)^{1 \over 2} $$

  $$ {\Delta E \over {1 \over 2\sqrt{2 \pi G}}} = {p \over c^3} + \rho $$

  $$ ds^2 = g_{\mu\nu}(x)dx^{\mu}dx^{\nu} + \phi^2(x)\kappa^2 A_{\mu\nu}(x)dx^{\mu})^2 $$
  $$ ds = (g_{\mu\nu}(x)dx^{\mu}dx^{\nu} + \phi^2(x)(\kappa^2 A_{\mu\nu}(x)dx^{\mu})^2)^{1 \over 2} $$

  $$ {\Delta E \over {1 \over 2\sqrt{2 \pi G}} - \rho} = c^3, {p \over 2 \pi} = c^3 $$
    

  $$ ds^2 = e^{-2kT(x)|\phi|}[\eta_{\mu\nu} + \bar{h}_{\mu\nu}(x)]dx^{\mu}dx^{\nu} + T^2(x) d\phi^2 $$

  $$ f_{z} = \int \begin{bmatrix} \sqrt{\begin{pmatrix} x_{1} & x_{2} & x_{3} \\
     y_{1} & y_{2} & y_{3} \end{pmatrix}
     \circ
   \begin{pmatrix} x_{1} & x_{2} & x_{3} \\
     y_{1} & y_{2} & y_{3} \end{pmatrix}}_{}
     \end{bmatrix}dxdydz $$

     $$ {\partial \over \partial f} \int (\sin2x)^2dx = ||x - y||^2 $$
	$$ ||\int [\nabla_{i}\nabla_{j}f]dm||^{{1 \over 2} + iy} = \mathrm{rot}(\mathrm{div},E,E_1) $$
	Maxwell of equation in fourth of power.
	$$ = 2 <f,h>, {V(x) \over f(x)} = \rho (x) $$
	$$ \int_M \rho (x)dx = \square \psi, - 2 <g, h> = \mathrm{div}(\mathrm{rot} E,E_1) $$
	$$ = - 2 R_{ij} $$
	Higgs field of space quality.
	$$ \mathcal{O}(x) = ||{\nabla_{i}\nabla_{j} \over S^2}\int [\nabla_{i}\nabla_{j}f\cdot g(x)dxdy]d\psi|| $$
	$$ = \int (\delta(x))^{2 \sin \theta \cos \theta} \log \sin \theta d \theta d\psi $$
	Helmander manifold of duality differential operator.
	$$ \delta \cdot \mathcal{O}(x) = [{||\nabla_{i}\nabla_{j} f d \eta|| \over 
	\int e^{2 \sin \theta \cos \theta} \cdot \log \sin \theta d \theta}] $$	
	Open set group in subset theorem, and helmandor manifold in singularity.
	$$ i\hbar \psi = ||\int {\partial \over \partial z} [{i(xy + \overline{yx}) \over z - \bar{z}}]dmd\psi|| $$
	Euler number. Complex curvature in Gamma function. Norm space.
	Wave of quantum equation in complex differential variable into integrate of rout.
	$$ \delta(x) \int_C R^{2 \sin \theta \cos \theta} = \Gamma(p + q) $$
	$$ (\square + m) \cdot \psi = (\nabla_{i}\nabla_{j}f|_{g_{\mu\nu(x)}} + v \nabla_{i}\nabla_{j}) $$
	$$ \int[m(x)(\mathrm{rot}\cdot\mathrm{div}(E,E_1))]dmd\psi $$
	Weak electric power and Maxwell equation combine with strong power.
	Private energy.
	$$ G_{\mu\nu} = \square \int \int \int (x,y,z)^3dxdydz $$
	$$ = {8 \pi G \over c^4}T^{\mu\nu} $$
	$$ {d \over dV}F = \delta(x) \int \nabla_{i}\nabla_{j}f d \eta_{\mu\nu} $$
	Dalanvelsian in duality of differential operator and global integrate variable operator.

	$$ x^n + y^n = z^n, \delta (x) \int z^n = {d \over dV}z^3, (x,y)\cdot(\delta^m,\partial^m) $$
	$$ = (x,y)\cdot(z^n,f) $$
	$$ n \bot x, n \bot y $$
	$$ = 0 $$
	Singularity of constance theorem.

	$$ \vee(\nabla_{i}\nabla_{j}f) \cdot \mathrm{XOR}(\square \psi)= {d \over df}\int_M F dV $$
Non-cognitive of summuate of equation and add manifold create with open set group theorem.
This group composite with prime of field theorem. Moreover this equations involved with symmetry dimension created.
	$$ \partial (x) \int z^3 = {d \over dV}z^3, \sum_{k=0}^{\infty}{1 \over (n + 1)^s} = \mathcal{O}(x) $$
	Singularity theorem and fermer equation concluded by this formula.
	$$ \int \mathcal{O}(x)dx = \delta(x)\pi(x)f(x) $$
	$$ \mathcal{O}(x) = \int \sigma(x)^{2\sin \theta \cos \theta}\log \sin \theta d\theta d\psi $$
	$$ y = x $$
	$$ \mathcal{O}(x) = ||{\nabla_{i}\nabla_{j}f \over S^2} \int [\nabla_{i}\nabla_{j}f \circ g(x)dxdy]|| $$
	$$ \lim\limits_{n \to 1}\sum_{k=0}^{\infty}{a_k \over a_{k+1}} = (\log \sin \theta dx)^{'} $$
	$$ = {\cos \theta \over \sin \theta} $$
	$$ = {x \over y} $$
	$$ \lim\limits_{n \to 1}\sum_{k=0}^{\infty}a_k f^k = {1 \over 1 - z} $$
	Duality of differential summuate with this gravity theorem.
	And Dalanverle equation involved with this equation.
	$$ \square \psi = {8\pi G \over c^4}T^{\mu\nu} $$
	$$ {\partial \over \partial f}\square \psi = 4 \pi G \rho $$
	$$ \int \rho (x) = \square \psi, {V(x) \over f(x)} = \rho(x) $$
	Dense of summuate stimulate with gravity theorem.

  今、大域的積分多様体に微分を入れると、大域的部分積分多様体が
   できる。ガンマ関数の大域的微分多様体と同型らしい。　


      大域的多重積分と大域的無限微分多様体を定義すると、
   Global assemble manifold defined with infinity of differential
   fields to determine of definition create from Euler product and
   Gamma function for Beta function being potten from integral manifold
   system. Therefore, this defined from global assemble manifold from
   being Gamma of global equation in topological expalanations.

   $$ \int f(x)dx = \int \Gamma(\gamma)^{'}dx_m $$
   $$ = 2 (\cos (i x \log x) - i \sin(i x \log x)) $$
   $$ \left({{\int f(x) dx} \over \log x}\right) = \lim_{\theta \to \infty}
    \left({{\int f(x) dx} \over \theta}\right) = 0,1 $$
    $$ e^{i \theta} = \cos \theta + i \sin \theta $$
    $$ \left({\int f(x) dx}\right)^{'} = 2 (i \sin(i x \log x) - \cos (i x \log
    x)) $$
    $$ = 2(- \cos(i x \log x) + i \sin(i x \log x)) $$
    $$ (\cos (i x \log x) - i \sin(i x \log x))^{'} $$
    $$ = {d \over {d {e^{i\theta}}}}\left((\cos, - \sin )\cdot
    (\sin, \cos)\right)  $$
    $$ \int \Gamma(\gamma)^{''}dx_m = \left({\int \Gamma(\gamma)^{'}}dx_m
    \right)^{\nabla L} = \left({{\int \Gamma dx_m} \cdot {{d \over d\gamma}
    \Gamma}}\right)^{\nabla L} \le 
    \left({\int \Gamma dx_m + {d \over d\gamma}\Gamma}\right)^{\nabla L} \le 
    \left({e^{f} - e^{-f}} \le {e^{-f} + e^{f}}\right)^{'}  $$
    $$ = 0,1 $$

    Then, the defined of global integral and differential manifold
    from being assembled world lines and partial equation of deprivate
    formula in Homology manifold and gamma function of global topology
    system, this system defined with Shwaltsshild cicle of norm from
    black hole of entropy exsented with result of monotonicity for
    constant of equations.
    以上 より、大域的微分多様体を大域的 2重微分多様体として、
    処理すると、ホモロジー多様体では、種数が１であり、特異点では、種数が０
    と計算されることになる。ガンマ関数の大域的微分多様体では、
    シュバルツシルト半径として計算されうるが、これを大域的2重微分で
    処理すると、ブラックホールの特異点としての解が無になる。




	Abel拡大 $ K/k $に対して、
	$$ f = {\pi}_{p}{f_p} $$
        類体論
	Artin記号を用いて、
       $$ \left({{\alpha, K/k} \over {p}}\right) 
       = \left({{K/k} \over {b}}\right) (\in G) $$
       $ {\alpha/ \alpha_0} \equiv 1 \pmod {f_p} $, 
      $ \alpha_0 \equiv 1 \pmod {f f_p^{-1}} \to  
      \alpha \in k $
      $ (\alpha_0) = {p}^{\alpha}b, \text{
	      $p$ と$b$は互いに素} $
      $b \to \text{相対判別式$\delta K/k$で互いに素}$
      この値は、補助数${\alpha_0}$の値の取り方によらずに、
      一意的に定まる。
      $$ \left({{\alpha, K/k} \over {p_{\infty}^{(j)}}}\right)  = 
      \text{1または0} $$

      これらをまとめた式が、Hilbertの剰余記号の判別式
      $$ \mathcal{\pi}_{p}\left({{\alpha,b} \over 
      {p}}\right) = 1 $$
      であり、
    　この式たちから、代数幾何の種数のノルム記号である、
$$ ||ds^2|| = \lim_{x \to \infty} [\delta(x) \int \int \int \pi \left( \sum_{k=0}^{\infty}
		{{}^n\sqrt{p},x \over n} \right)
^{1 \over 2}d \tau]^{\mu\nu} $$
     が求まり、
     $$ p^{\alpha} n = {}^n \sqrt{p} $$
     $$ n^{{}^n\sqrt{p}} = \bigoplus{(i \hbar^{{{\nabla)}}^{\oplus 
     L}}} $$
     $$ = n^{p^{{1 \over n}}} = n^{{-n}^{p}}  $$
     $$ = \int \Gamma(\gamma)^{'}dx_m = e^{-x \log x} $$
     となり、
     kの素イデアルの密度Mに対して、
     $$ \lim_{s \to {1 + 0}} \sum_{{p \in M}} {
	     1 \over (N(p))^{s}}/ \log {1 \over 
     {s - 1} } $$
     $$ = \text{Mの密度(density)} $$

     $$ \alpha(f{d \over dt},g{d \over dt}) = \int_s \begin{vmatrix}
	     f^{'} & f^{''} \\
	     g^{'} & g^{''} \end{vmatrix}dt,
	     \mathcal{B}(f{d \over dt},g{d \over dt},h{d \over dt}) =
	     \int_s \begin{vmatrix} f & f^{'} & f^{''} \\
		     g & g^{'} & g^{''} \\
	     h & h^{'} & h^{''} \end{vmatrix}dt $$
    これらは、Gul'faid-Fuksコホモロジーの概念を局所化することにより、
    形式的ベクトル和のつくるコホモロジーとして、導かれている。
$$ || \int f(x) ||^2 \to \int \pi r^2 dx \cong V(\tau) $$

 This equation is part of summulate on entropy formula.

 $$ {\int |f(x)|^2 \over f(x)} dx $$
$$V(\tau) \to mesh $$

 $$ \int \begin{vmatrix} x_1 & x_2 & x_3 \\
    y_1 & y_2 & y_3 \\
    z_1 & z_2 & z_3 \end{vmatrix} dv({\tau}) $$
 $$ \begin{vmatrix} x_1 & x_2 & x_3 \\
    y_1 & y_2 & y_3 \\
    z_1 & z_2 & z_3 \end{vmatrix}_{dx=v}  $$


   $$ z_y = a_x + b_y + c_z $$
$$ dz_y = d(z_y) $$

$$ [f, f^{-1}] = ff^{-1} - f^{-1}f $$

  Universe of extend space has non-integrate rout entropy in sufficient of mass in darkmatter,
this reason belong for cohomology and heat equation concerned. This expanded mass match is darkmatter
of extend. And these heat area consume is darkmatter begun to start big-ban system invite of universe.

$$ V(\tau) = \int \tau(q)^{-{n \over 2}} \mathrm{exp}(-{1 \over \sqrt{2 \tau(q)}}L(x)dx) + O(N^{-1}) $$
$$ {1 \over \tau}({N \over 2} + \tau(2 \Delta f - |\nabla f|^2 + R) + f)\mathrm{mod}N^{-1} $$
	$$ \Delta E = -2(T - t)|R_{ij} + \nabla_{i} \nabla_{j} f - {1 \over 2(T - t)}g_{ij}|^2 $$
	$$ {d \over df}F = {2\int (R + \nabla_{i} \nabla_{j} f)^2 \over - (R + \Delta f)}dm $$
	$$ {d \over dt}g_{ij}(t) = - 2 R_{ij} $$
	$$ dx = (g_{\mu\nu}(x)^2dx^2 - g_{\mu\nu}(x)dxg_{\mu\nu})^{1 \over 2} $$
	$$ ds^2 = - N(r)^2dt^2 + \psi^2(r)(dr^2 + r^2d\theta^2) $$
  $$ f_{z} = \int \begin{bmatrix} \sqrt{\begin{pmatrix} x_{1} & x_{2} & x_{3} \\
     y_{1} & y_{2} & y_{3} \end{pmatrix}
     \circ
   \begin{pmatrix} x_{1} & x_{2} & x_{3} \\
     y_{1} & y_{2} & y_{3} \end{pmatrix}}_{}
     \end{bmatrix}dxdydz $$

     $$ \sum^{\infty}_{n=0}a_1x^1 + a_2x^2 \dots a_{n-1}x^{n-1} \to \sum^{\infty}_{n=0} a_nx^n \to \alpha $$

$$ f = n\nu\lambda , \lambda = {x \over l},
 \int dn \nu \lambda = f(x), 
 xf(x) = F(x),
 [f(x)] = \nu h$$


    石岡先生のところで見て、石川先生と河相先生、脇本先生のところで、
    確かめている。

    $$ {2 e^{x_n t} \over {e^{t} + 1}} = \sum_{n=0}^{\infty} E_n(x){{t^x} 
    \over {n!}} $$
    $$ E_n(x) = \sum_{k=0}^{n} \begin{pmatrix} n \\
    k\end{pmatrix}  a_n x^{n -k} $$

	   $$ y =  1 + \tan^2 x = {1 \over {\cos x}} $$
	   $$ y^{'} = {{\cos^{'}x} \over {\cos^2 x}} $$
           $$ -{{\cos x} \over (\cos x)^{'}}(\sin x)^{'} = z_n $$
	   $$ {1 \over y^{'}} = z^2_n$$
           $$ z^n = - 2 e^{x \log x} $$

	   ガンマ関数の大域的積分多様体は、結局は、フェルマーの定理
	   だった。

         代数的計算手法のために$ \oplus L $を使っている。
そのために、冪乗計算と商代数の計算が、乗算で楽に見えるようになっている。
            微分幾何の量子化は、代数幾何の量子化の計算になっている。
  加減乗除が初等幾何であり、大域的微分と積分が、現代数学の代数計算の
  簡易での楽になる計算になっている。
  初等代数の計算は、
  $$ \bigoplus({i \hbar^{\nabla}})^{\oplus L} $$
   $$m \bigoplus({i \hbar^{\nabla}})^{\oplus L} +  
   n\bigoplus({i \hbar^{\nabla}})^{\oplus L} = 
   (m+n) \bigoplus({i \hbar^{\nabla}})^{\oplus L}   $$
   $$ m\bigoplus({i \hbar^{\nabla}})^{\oplus L} -  
   n\bigoplus({i \hbar^{\nabla}})^{\oplus L} = 
   (m-n) \bigoplus({i \hbar^{\nabla}})^{\oplus L}   $$
   $$ \bigoplus({i \hbar^{\nabla}})^{{\oplus L}^m} \times 
      \bigoplus({i \hbar^{\nabla}})^{{\oplus L}^n} =
   \bigoplus({i \hbar^{\nabla}})^{\oplus L^{m+n}} $$
  $$ {{\bigoplus({i \hbar^{\nabla}})^{\oplus L^m}}  \over  
  {\bigoplus({i \hbar^{\nabla}})^{\oplus L^n}}} =  
  \bigoplus({i \hbar^{\nabla}})^{\oplus L^{m\over n}} $$
  大域的計算での微分と積分は、
  $$ \left(\bigoplus({i \hbar^{\nabla}})^{\oplus L}\right)^{df}
  =  {\left(\bigoplus({i \hbar^{\nabla}})^{\oplus L}\right)^{
	  \bigoplus({i \hbar^{\nabla}})^{\oplus L}}}^{'}  $$
  $$ = \bigoplus({i \hbar^{\nabla}})^{\oplus L^{'}} $$
   $$ \int \bigoplus({i \hbar^{\nabla}})^{\oplus L} dx_m $$
   $$ \bigoplus({i \hbar^{\nabla}})^{\oplus L} = n $$
   大域的積分の代数幾何の量子化での計算は、ガンマ関数になっている。
   $$ {n^{L+1} \over {n+1}} = \int e^x x^{1 - t} dx $$
   代数幾何の量子化の因数分解は、加減乗除の式のm,nの
   組み合わせ多様体でのガンマ関数同士の計算になる。


       This estern with Gamma function resteamed from being riging to Beta function in Thurston
   Perelman manifold. This field call all of theorem to architect with Space ideal of quantum level.
     This theorem will be estern the man to be birth with Japanese person.
   This person pray with be birth of my son.
   This pray call work to be being name to say me pray.
   This pray resteam me to masterbation and this play realized me Gakkari.
   Aya san kill me to be played.

     I like this poem to proof with English moreover Japanese language loved from me.
     And this crystal proof released me to write English and Japanese language to discover them
     from mathmatics theorems.

   $$ \left(\bigoplus({i \hbar^{\nabla}})^{\oplus L}+m\right) 
   \left(\bigoplus({i \hbar^{\nabla}})^{\oplus L}+n\right) $$
  $$ = {n^{L+1} \over {n+1}} = t \int (1 - x)^n x^{m - 1} dx $$



     This equation esterminate with Beta function in Gamma function riginged from telphone
  to world line surface. And this ringed have with Algebra manifold of differential geometry in
  quantum level. This write in English language. Moreover that' cat call them to birth of Japanese cats.
  And moreover, I birth to name with Japanese Person. 
  And, this theorem certicefate the man to birth Diths Person.
  This stimeat with our constrate with non relate person and cat.


  $$ = \int x^{m-1}(1 - x)^{n-1} dx $$


  $$ \nabla ({i \hbar^{\nabla}})^{\oplus L}, \bigoplus ({i \hbar^{\nabla}})^{
	   \oplus L}, \square ({i \hbar^{\nabla}})^{\oplus L} $$
   $$ \boxtimes (i \hbar^{\nabla})|_{dx_m}^L, \boxplus (i \hbar^{\nabla}|_
   {dx_m}^L) $$



  $$ x^{{1 \over 2} + iy} = e^{x \log x} $$
  それゆえに、この定数はゼータ関数と微分幾何の量子化を因数にもつ
  素因子分解の式にもなっている。
  Therefore, this product also constructed with differential geometry
  of quantum level equation and zeta function.
  $$ = C $$ 
  そして、この関数はオイラーの定数から広中平祐定理による4重帰納法の
  オイラーの公式からの多様体積分へとのサーストン空間のスペクトラム関数
  ともなっている。
 
  $$ (\psi^{out}(h_1, \dotsb h_n), \psi^{out}(h^{'}_1, \dotsb h^{'}_{{n^{'}}}
  )) $$
  $$ = \delta_{n n^{'}} \sum_{p} \Pi^{n}_{j=1} (h,h_{p}(
  j)) $$
  このHaag-Ruelleの散乱理論からのFucks空間の$ \psi^{out} $の１粒子状態が
  $\sum^{\oplus}R_1(m)$上の$\mathcal{K^{0}}$への閉部分空間への
  $R^{out}$のユニタリ写像として、複素空間の種数の重ね合わせをノルム空間と
  見なして、これが代数幾何の量子化と同型と言えることを、
  これらの式たちから示しているのが、Hilbert空間のノルムである。　

\end{document}

